%==============================================================================
\section{Unsupervised Learning --- Học máy không giám sát}
\label{sec:unsupervised-learning}
%==============================================================================

\subsection{Unsupervised Machine Learning là gì?}

\begin{dinhnghia}[title={Học không giám sát}]
Unsupervised learning học pattern và cấu trúc trong dữ liệu \textbf{không} cần giám sát con người.
Thuật toán phân tích dữ liệu để khám phá pattern ẩn trong tập \textbf{không nhãn}.
\end{dinhnghia}

\begin{vidu}[title={Minh hoạ unsupervised learning}]
\tsimg{11_mo_hinh_va_cac_kieu_hoc/img03.jpg}
\end{vidu}

\begin{ynghia}[title={Khác supervised}]
Không có quyền ``thoải mái'' như supervised với dữ liệu đã gán nhãn sẵn; cần trích pattern hữu ích từ đầu vào.
\end{ynghia}

\begin{vidu}[title={Tóm tắt một câu}]
\begin{itemize}
  \item Là cách tiếp cận/kỹ thuật ML
  \item Dùng thuật toán ML
  \item Tìm pattern/cấu trúc ẩn
  \item Trong dữ liệu không có giám sát con người
\end{itemize}
\end{vidu}

\begin{phuongphap}[title={Một số hướng và ví dụ}]
Association, clustering, dimensionality reduction.
Ví dụ thuật toán: K-means clustering, K-nearest neighbors, \ldots
\end{phuongphap}

\begin{vidu}[title={Ví dụ trực giác: dữ liệu cử tri}]
Thay vì huấn luyện máy dự đoán $Y$ cho mỗi $X$ như supervised, ta hỏi:
``Với tập dữ liệu khổng lồ $X$, bạn cho biết gì về $X$?''
Hoặc: ``Có thể chia $X$ thành năm nhóm tốt nhất nào?''
Hoặc: ``Ba feature nào hay xuất hiện cùng nhau nhất trong $X$?''
--- Đó là tinh thần của unsupervised learning.
\end{vidu}

\subsection{Cách Unsupervised Learning hoạt động}

\begin{phuongphap}[title={Quy trình huấn luyện}]
\begin{itemize}
  \item Thuật toán tự học (self-learning) huấn luyện trên dữ liệu \textbf{không nhãn}.
  \item Tùy tác vụ (clustering, association, \ldots) chọn thuật toán phù hợp.
  \item Thuật toán tự suy ra quy tắc dựa trên độ tương đồng, pattern, khác biệt --- không có nhãn hay pre-training.
  \item Kết quả: \textbf{mô hình unsupervised}, sẵn sàng cho clustering, association hoặc giảm chiều.
\end{itemize}
\end{phuongphap}

\begin{ynghia}[title={Phù hợp bài toán phức tạp}]
Mô hình unsupervised phù hợp tổ chức dataset lớn thành các cụm.
\end{ynghia}

\subsection{Ba nhóm phương pháp}

\subsubsection{Clustering}

\begin{dinhnghia}[title={Clustering}]
Gom tập đối tượng/điểm dữ liệu thành cụm theo độ tương đồng.
Điểm trong cùng cụm phải \textbf{giống nhau hơn} điểm ở cụm khác.
Đôi khi gọi là \textbf{unsupervised classification} vì kết quả giống phân loại nhưng không có lớp định sẵn.
\end{dinhnghia}

\begin{ynghia}[title={Một số thuật toán}]
\textbf{K-Means}: gán điểm vào một trong $K$ cụm theo khoảng cách tới tâm; cập nhật centroid lặp đến khi ổn định.

\textbf{Mean Shift}: tìm vùng mật độ dữ liệu cao; dịch mean của từng điểm về vùng dày nhất.

\textbf{Gaussian Mixture Models}: mô hình xác suất kết hợp nhiều phân phối Gaussian để xác định điểm thuộc thành phần nào.
\end{ynghia}

\subsubsection{Association Rule Mining}

\begin{ynghia}[title={Association}]
Kỹ thuật dựa trên luật để tìm liên kết giữa tham số trong dataset lớn.
Phổ biến trong Market Basket Analysis; hỗ trợ quyết định và recommendation engine.
Thuật toán chính: \textbf{Apriori} --- tìm item lặp lại thường xuyên và luật kết hợp trong dữ liệu giao dịch.
\end{ynghia}

\subsubsection{Dimensionality Reduction}

\begin{ynghia}[title={Giảm chiều}]
Chọn tập feature đại diện để giảm số chiều mỗi mẫu --- tránh curse of dimensionality khi có hàng triệu feature.
Thuật toán phổ biến: Principal Component Analysis (PCA), Missing Value Ratio, SVD, Autoencoders.
\end{ynghia}

\subsection{Thuật toán Unsupervised Learning}

\begin{phuongphap}[title={Danh sách thường dùng}]
\textbf{Clustering}: K-Means, Hierarchical, Mean-shift, DBSCAN, HDBSCAN, BIRCH, Affinity Propagation, Agglomerative.

\textbf{Association}: Apriori, Eclat, FP-growth.

\textbf{Giảm chiều}: PCA, Autoencoders, SVD.
\end{phuongphap}

\subsection{Ưu điểm}

\begin{tinhchat}[title={Ưu điểm}]
\begin{itemize}
  \item Không cần dữ liệu có nhãn --- dễ và rẻ hơn.
  \item Khám phá pattern/quan hệ ẩn $\Rightarrow$ insight và ra quyết định.
  \item Phù hợp clustering, anomaly detection, giảm chiều.
\end{itemize}
\end{tinhchat}

\subsection{Nhược điểm}

\begin{luuy}[title={Nhược điểm}]
\begin{itemize}
  \item \textbf{Khó đánh giá}: không có nhãn chuẩn để đo hiệu năng.
  \item \textbf{Kết quả có thể kém chính xác}: đặc biệt khi dữ liệu nhiễu hoặc thuật toán không biết đầu ra đúng.
\end{itemize}
\end{luuy}

\subsection{Ứng dụng}

\begin{phuongphap}[title={Thực tế}]
\begin{itemize}
  \item \textbf{Customer Segmentation}: gom khách theo mua hàng, hoạt động, sở thích.
  \item \textbf{Anomaly Detection}: phát hiện pattern bất thường (gian lận, an ninh mạng).
  \item \textbf{Recommendation Engines}: phân tích hành vi khách hàng, marketing cá nhân hoá.
  \item \textbf{NLP}: ví dụ Google phân loại bài báo tin tức.
\end{itemize}
\end{phuongphap}

\subsection{Anomaly Detection}

\begin{dinhnghia}[title={Phát hiện bất thường}]
Tìm sự kiện/quan sát hiếm, hiếm khi xảy ra.
Dùng kiến thức đã học để phân biệt điểm bất thường và bình thường.
Clustering, KNN có thể hỗ trợ dựa trên feature.
\end{dinhnghia}
